\documentclass[11pt, a4paper]{article}

\usepackage[T1]{fontenc}
\usepackage[utf8]{inputenc}
\usepackage[spanish]{babel}
\usepackage[a4paper, margin=2.5cm]{geometry}
\usepackage{amsmath, amssymb}
\usepackage{booktabs}
\usepackage{enumitem}
\usepackage{listings}
\usepackage{xcolor}
\usepackage{hyperref}

\definecolor{codigobg}{RGB}{245,245,245}
\definecolor{keyword}{RGB}{0,100,150}
\lstset{language=R, basicstyle=\ttfamily\small, backgroundcolor=\color{codigobg},
  keywordstyle=\color{keyword}\bfseries, frame=single, breaklines=true, showstringspaces=false}

\setlength{\parindent}{0pt}
\pagestyle{plain}

\begin{document}

\begin{center}
{\large\bfseries PONTIFICIA UNIVERSIDAD JAVERIANA}\\[3pt]
Facultad de Ciencias Econ\'omicas y Administrativas\\[2pt]
Estad\'istica 2026-II --- Prof.\ Jaime Polanco Jim\'enez\\[6pt]
{\LARGE\bfseries Problem Set 6}\\[3pt]
{\large ANOVA, Pruebas No Param\'etricas y Asociaci\'on Categ\'orica}\\[6pt]
Asignado: Septiembre 25 (Sesi\'on 6) \quad$|$\quad Entrega: Octubre 2 (Sesi\'on 7)
\end{center}

\vspace{6pt}
\rule{\linewidth}{0.5pt}
\vspace{2pt}
{\centering\small Repositorio: \url{https://github.com/polanco-jaime/Curso-Estadistica-2026-3}\par}
\vspace{6pt}

\textbf{Instrucciones.} Este Problem Set es individual y se resuelve en R sobre
microdatos ICFES. La entrega es digital al inicio de la clase indicada.
El c\'odigo R debe incluirse.

\vspace{10pt}

\section*{Enunciado}

En este Problem Set explorar\'a diferencias en el puntaje de ingl\'es de
estudiantes de Negocios Internacionales (NI) seg\'un diferentes variables
categ\'oricas, usando ANOVA, pruebas no param\'etricas y medidas de asociaci\'on.

\begin{enumerate}[label=\arabic*)]
  \item \textbf{ANOVA de una v\'ia por estrato socioecon\'omico:} ¿var\'ia el
        puntaje de ingl\'es en Saber Pro (\texttt{MOD\_INGLES\_PUNT}) seg\'un el
        estrato socioecon\'omico de los estudiantes de NI?
        \begin{enumerate}[label=\alph*)]
          \item Plantee las hip\'otesis del ANOVA:
                \[
                  H_0: \mu_1 = \mu_2 = \cdots = \mu_k \quad \text{vs.} \quad
                  H_1: \text{al menos una media difiere}
                \]
                donde $k$ es el n\'umero de estratos.
          \item Verifique los supuestos del ANOVA:
                \begin{itemize}
                  \item Normalidad por grupo (prueba de Shapiro-Wilk o Q-Q plots)
                  \item Homocedasticidad (prueba de Levene)
                \end{itemize}
          \item Aplique ANOVA usando \texttt{aov()} y reporte la tabla ANOVA
                (fuentes de variaci\'on, sumas de cuadrados, $F$, $p$-valor).
          \item Si rechaza $H_0$: ¿cu\'al es el porcentaje de variabilidad
                explicado por el estrato ($\eta^2$)?
        \end{enumerate}

  \item \textbf{Comparaciones m\'ultiples con Tukey HSD:} si el ANOVA fue
        significativo, identifique cu\'ales pares de estratos difieren.
        \begin{enumerate}[label=\alph*)]
          \item Aplique la prueba de Tukey HSD (\texttt{TukeyHSD()}).
          \item Reporte los intervalos de confianza ajustados al 95\% para
                todas las comparaciones por pares.
          \item ¿Cu\'ales pares de estratos muestran diferencias significativas?
          \item Grafique los intervalos de confianza de Tukey.
        \end{enumerate}

  \item \textbf{Prueba de Mann-Whitney por g\'enero:} ¿difiere el puntaje de
        ingl\'es entre hombres y mujeres en NI?
        \begin{enumerate}[label=\alph*)]
          \item Justifique por qu\'e puede ser apropiado usar una prueba no
                param\'etrica (Mann-Whitney) en lugar de la prueba $t$.
          \item Plantee $H_0$ y $H_1$ en t\'erminos de medianas o rangos.
          \item Aplique \texttt{wilcox.test()} y reporte el estad\'istico $W$
                y el $p$-valor.
          \item Interprete el resultado: ¿rechaza $H_0$ al 5\%? ¿Qu\'e grupo
                tiene mayor puntaje mediano?
        \end{enumerate}

  \item \textbf{Prueba $\chi^2$ de independencia: g\'enero $\times$ nivel MCER:}
        ¿existe asociaci\'on entre el g\'enero y el nivel de ingl\'es seg\'un
        el MCER (A$-$, A1, A2, B1, B2) en NI?
        \begin{enumerate}[label=\alph*)]
          \item Construya una tabla de contingencia cruzando g\'enero con
                nivel MCER (\texttt{table()}).
          \item Plantee $H_0$ (independencia) y $H_1$ (asociaci\'on).
          \item Aplique la prueba $\chi^2$ de independencia
                (\texttt{chisq.test()}) y reporte el estad\'istico $\chi^2$,
                grados de libertad y $p$-valor.
          \item Verifique el supuesto: ¿todas las frecuencias esperadas son
                $\geq 5$?
        \end{enumerate}

  \item \textbf{Tama\~no del efecto: $V$ de Cram\'er:} cuantifique la fuerza
        de la asociaci\'on g\'enero--MCER.
        \begin{enumerate}[label=\alph*)]
          \item Calcule la $V$ de Cram\'er:
                \[
                  V = \sqrt{\frac{\chi^2}{n \cdot (k - 1)}}
                \]
                donde $k = \min(\text{filas}, \text{columnas})$.
          \item Interprete $V$ seg\'un la convenci\'on (d\'ebil: $V < 0.1$;
                moderada: $0.1 \leq V < 0.3$; fuerte: $V \geq 0.3$).
          \item ¿La asociaci\'on encontrada (si es significativa) tiene
                relevancia pr\'actica?
        \end{enumerate}

  \item \textbf{Conclusiones:} sintetice los hallazgos:
        \begin{enumerate}[label=\alph*)]
          \item ¿El estrato socioecon\'omico explica diferencias en el puntaje
                de ingl\'es? ¿Cu\'ales estratos difieren?
          \item ¿Existe brecha de g\'enero en ingl\'es en NI?
          \item ¿El nivel de ingl\'es MCER est\'a asociado con el g\'enero?
        \end{enumerate}
\end{enumerate}

\vspace{10pt}
\rule{\linewidth}{0.4pt}
\begin{center}
{\small Cada entrega completa suma $+0.0625$ a la nota final.}
\end{center}

\end{document}
